In this section, we introduce the verifiable reward modeling setup in the RLVR framework and the concept of ``master key'' attacks.
% followed by our approach to training a robust reward model to defend against them.

\paragraph{Verifiable Reward Model.}
Reinforcement Learning with Verifiable Rewards (RLVR)~\citep{luong2024reft,lambert2024t,guo2025deepseek,su2025crossing} focuses on a reference-based setting, where the reward signal is provided by either a rule-based function or a generative, LLM-based judge. At each step of RLVR training, the reward model receives a question $q$, a response $o$ generated by the policy model, and a reference answer $a^*$, and produces a binary signal $y \in \{\texttt{YES}, \texttt{NO}\}$ that determines whether $o$ aligns with $a^*$ given $ q$. 


Formally, the LLM judge defines a function:
\[
J(q, a^*, o) \rightarrow \{\texttt{YES}, \texttt{NO}\}
\]
 This judgment translates directly into a reward signal, which guides the training of the policy model: a positive reward ($R=1$) for a \texttt{YES} and a zero reward ($R=0$) for a \texttt{NO}.
Thus, the accuracy and reliability of this judgment directly affect the policy model’s training. Any systematic failures or false positive rewards in the verification process can mislead the learning trajectory.

\paragraph{Master Keys.}
In this work, we identify a family of adversarial patterns, termed \emph{``master keys''}. When used as responses, these patterns can \emph{surprisingly} trigger false positive judgments from a wide range of LLM judges, even though they are semantically meaningless for solving the task. This effect holds across diverse $(q, a^*)$ from various data domains.
These patterns can be divided into two categories: (1) \textbf{Non-word symbols} including punctuation such as \emph{``.''}, \emph{``:''} and (2) \textbf{Reasoning openers} which involve natural language expressions that signal the start or structure of a reasoning process, but do not yet contribute substantive content (e.g., \emph{``Thought process:''}, \emph{``Solution''}, \emph{``Let’s solve this problem step by step.''}).

 

Despite offering little meaningful contribution to problem-solving, these expressions are often accepted as correct by multiple LLM judges across diverse datasets. We show that such false positive rewards persist even with model-specific evaluation prompts and with state-of-the-art LLMs, including GPT-4o, Claude-4, Qwen2.5-72B-Instruct, as well as specialized reference-based generative reward models, including Qwen2.5-7B-Instruct-RLVR~\citep{su2025crossing}\footnote{Throughout this work, we shall refer to this model as \emph{Multi-sub RM} for simplicity.} and Omni-Judge~\citep{gao2024omnimathuniversalolympiadlevel}. This reveals a critical and underexplored vulnerability in the core mechanics of reward modeling: the verifier, designed to filter out invalid or incorrect answers, can be manipulated by trivial, superficial content, resulting in false positives. This undermines the integrity of any pipelines (e.g., RLVR) that rely on generative verifiers for feedback.




